%% ma: L12-B08-0014
%% lop: 12
%% chuong: 2
%% bai: 8
%% dang: 12.8.7
%% loai: TN4
%% muc: NB
%% dap_an: "A"
%% nguon: "(NBV)-ĐỀ SỐ 8-MỖI NGÀY 1 ĐỀ THI 2025.pdf (D08, MỖI NGÀY 1 ĐỀ - Đề số 8), Phần I Câu 9"
%% kien_thuc: "Bài 8 (biểu thức tọa độ, tọa độ trung điểm)"
%% trang_thai: chinh-thuc
%% ghi_chu: "Đáp án gốc A đúng. Cùng công thức trung điểm với L12-B08-0008 (dạng 12.8.7) nhưng hỏi ngược (tìm đầu mút)."
\begin{ex}
Trong không gian $Oxyz$, cho hai điểm $A(3;1;-2)$ và $B(-1;3;1)$. Biết rằng $B$ là trung điểm của đoạn thẳng $AC$. Toạ độ của điểm $C$ là
\choice{\True $(-5;5;4)$}{$\left(1;2;-\dfrac12\right)$}{$(7;-1;-5)$}{$\left(2;-1;-\dfrac32\right)$}
\loigiai{$B$ là trung điểm của $AC$ nên $x_C=2x_B-x_A=-5$, $y_C=2y_B-y_A=5$, $z_C=2z_B-z_A=4$. Vậy $C(-5;5;4)$.}
\end{ex}
